% generated by experiments/phase2/paper_evidence.py; do not edit
\begin{table*}[t]
\centering
\small
\caption{Worked example of Eq.~\ref{eq:crossover} at width 2048, population 256 and eight devices; all times in milliseconds. $C_A$ and $C_B$ are the separately measured reconstruction times. $T_{\mathrm{add}}$ is the time the all-reduce adds to the split reconstruction, measured in context; it stands in for $T_{\mathrm{ar}}(4P)$. The expected difference is $C_A-C_B+T_{\mathrm{ag}}-T_{\mathrm{add}}$, where $T_{\mathrm{ag}}$ is replication's second coefficient gather. Positive differences favor splitting. $t_A$ is the replicated generation time, the denominator of the ratio in Figure~\ref{fig:f2}.}
\label{tab:worked}
\begin{tabular}{llrrrrrr}
\toprule
platform & variant & $C_A$ & $C_B$ & $T_{\mathrm{add}}$ & expected $t_A-t_B$ & measured $t_A-t_B$ & $t_A$ \\
\midrule
A100 & seed & 61.05 & 7.78 & 1.25 & 52.03 & 54.30 & 113.26 \\
A100 & rank 1 & 0.77 & 0.35 & 1.62 & -1.19 & -1.40 & 19.88 \\
v5e & seed & 233.29 & 29.51 & 2.03 & 201.75 & 213.30 & 467.58 \\
v5e & rank 1 & 0.55 & 0.33 & 2.36 & -2.14 & -2.04 & 9.73 \\
\bottomrule
\end{tabular}
\end{table*}
